%قبل از شروع مسئله ترجیحاً چیزی ننویسید. اما اگر خواستید بنویسید، در قالب لیستی (فقرات و شمارش) بنویسید.
%برای اضافه کردن نمره، بعدی از شروع‌مسئله، [x] را اضافه کنید و به جای x  نمره مورد نظر را قرار دهید. برای زیرمسئله‌ها هم به همین صورت است.
%\شروع‌مسئله[10]
\شروع‌مسئله [9]
با مرتب‌سازی سریع در کلاس آشنا شدید. پیچیدگی زمانی این الگوریتم‌ را در حالات زیر بررسی و هر یک را اثبات کنید. قبل از پاسخ به این سؤال، نوع مرتب‌سازی سریع را (از سه نوع گفته شده در کلاس) به دلخواه انتخاب کنید و بر اساس آن به سؤالات پاسخ دهید. در پاسخ خود به صراحت اشاره کنید که کدام نوع مرتب‌سازی سریع را انتخاب کرده‌اید.

\پاسخ{	در این پاسخ نسخه‌ای از مرتب‌سازی سریع (غیرتصادفی) از دکتر قدسی در نظر گرفته شده است. در این نسخه عضر اول به عنوان لولا (محور) درنظر گرفته می‌شود. همچنین آرایه به دو قسمت (و نه سه قسمت) تقسیم می‌شود و این تقسیم به گونه‌ای پیاده‌سازی شده است که در صورتی که عناصر آرایه برابر باشند، آرایه به دو قسمت تقریباً مساوی تقسیم می‌شود.
	
}

\شروع{زیرمسائل}
 
\زیرمسئله[3] آرایۀ ورودی مرتب باشد.

\پاسخ{
		$O(n^2)$.
	هر تقسیم	در 
	$O(n)$
	انجام می شود و عنصر محور هر دفعه کوچکترین عنصر است. پس در کل 
	$O(n^2)$
	طول می کشد تا اجرای الگوریتم تمام شود.
	
}

\زیرمسئله[3] همۀ اعضای آرایۀ ورودی برابر باشند.

\پاسخ{
			$O(n\log n)$.
	هر 
	تقسیم
	در 
	$O(n)$
	انجام می شود، اما محور هر دفعه به وسط زیر آرایۀ مورد نظر می‌رود و آرایه را به دو بخش مساوی تقسیم می‌کند.
	
}

\زیرمسئله[3] آرایۀ ورودی به صورت نزولی مرتب شده باشد.

\پاسخ{
		    $O(n^2)$.
	آرایه در هر مرحله به دو قسمت به اندازه های 1 و 
	$n-1$
	تقسیم می شود و چون زمان اجرای تقسیم
	$O(n)$
	است، زمان اجرای الگوریتم در کل برابر است با 
	$O(n^2)$.
	
}

\پایان{زیرمسائل}

%اگر مسئله چند قسمت داشت به این صورت عمل کنید:
%\شروع{زیرمسائل}
%\زیرمسئله[0] % مسئلۀ ۱آ
%\زیرمسئله % مسئلۀ ۱ب
%\زیرمسئله % مسئلۀ ۱پ
%\زیرمسئله % مسئلۀ ۱ت
%\پایان{زیرمسائل}




%%%%%%%%%%%%%%%%%%%%%%%%%%%%%%%%%%%%%%%%%%%%%%%%%%%%%
% برخی دستورات مفید برای تمرین‌ها. %
%%%%%%%%%%%%%%%%%%%%%%%%%%%%%%%%%%%%%%%%%%%%%%%%%%%%%

% فرمول‌های نمونه:
%$f(x) = \Theta(x)$
%$T(x, y) \leq \log(x) + 2^y + \binom{2n}{n}$
% {\tt code\_function}


% لیست موارد بدون شماره‌گذاری:
%\شروع{فقرات}
%    \فقره مورد اول 
%        \شروع{فقرات}
%            \فقره مورد اول داخلی 
%                \شروع{فقرات}
%                    \فقره اولین مورد خیلی داخلی 
%                    \فقره دومین مورد خیلی داخلی 
%                \پایان{فقرات}
%            \فقره مورد دوم داخلی 
%        \پایان{فقرات}
%    \فقره مورد دوم 
%\پایان{فقرات}

% لیست موارد شماره‌گذاری شده:
%\شروع{شمارش}
%    \فقره یک مورد 
%    \فقره یک مورد دیگر 
%    \فقره باز هم یک مورد دیگر
%\پایان{شمارش}

% برای نوشتن یک فرمول چند خطی که یک شماره بخوره. برای حالت‌های دیگر به سایت زیر مراجعه کنید:
% https://www.overleaf.com/learn/latex/Aligning_equations_with_amsmath
%\begin{align} 
%    \begin{split}
%        (x+y)^3 &= (x+y)^2(x+y) \\
%                &= (x^2+2xy+y^2)(x+y) \\
%                &= (x^3+2x^2y+xy^2) + (x^2y+2xy^2+y^3) \\
%                &= x^3+3x^2y+3xy^2+y^3
%    \end{split}                                 
%\end{align}

% نوشتن ماتریس و آرایه دوبعدی
%\begin{align}
%    A = 
%    \begin{bmatrix}
%        A_{11} & A_{21} \\
%        A_{21} & A_{22}
%    \end{bmatrix}
%\end{align}

% اضافه کردن تصویر با شرح. برای اضافه کردن تصویر بدون شرح، فقط دستور درج‌تصویر را استفاده کنید.
%\شروع{شکل}[h!]
%\تنظیم‌ازوسط
%  \درج‌تصویر[پهنا=0.5\پهنای‌متن]{img.jpg}
%\شرح{شرح تصویر}
%\برچسب{تصویر:۱}
%\پایان{شکل}
%\پایان{زیرمسائل}

%کد زیر را ببینید.
%%\lstinputlisting{source_filename.py}
%\begin{latin}
%\begin{lstlisting}
%# Import the modules
%import sys
%import random
%
%ans = True
%
%while ans:
%question = raw_input("Ask the magic 8 ball a question: (press enter to quit) ")
%
%answers = random.randint(1,8)
%
%if question == "":
%sys.exit()
%
%elif answers == 1:
%print "It is certain"
%
%elif answers == 2:
%print "Outlook good"
%
%elif answers == 3:
%print "You may rely on it"
%
%elif answers == 4:
%print "Ask again later"
%
%elif answers == 5:
%print "Concentrate and ask again"
%
%elif answers == 6:
%print "Reply hazy, try again"
%
%elif answers == 7:
%print "My reply is no"
%
%elif answers == 8:
%print "My sources say no"
%\end{lstlisting}
%\end{latin}
